\documentclass[11pt,letterpaper]{article}
\usepackage{inputenc}
\usepackage{amsmath}
\usepackage{amsfonts}
\usepackage{amssymb}
\usepackage{graphicx}
\usepackage{enumitem}
\usepackage{fancyhdr}
\usepackage{tikz}
\usepackage{hyperref}
\usepackage[T1]{fontenc}
\usepackage{graphicx}
\usepackage{color}
\usepackage{bm}
\usepackage{ mathrsfs }
\usepackage{multirow}

\usepackage{algorithm}
\usepackage[noend]{algpseudocode}
\usepackage[backend=bibtex]{biblatex}
\definecolor{lightgray}{gray}{0.5}
\setlength{\parindent}{0pt}
 \usepackage[top=3cm, bottom=2cm, left=1.9cm, right=1.9cm]{geometry}

% \usepackage[table,xcdraw]{xcolor}
\usepackage{xcolor}
\colorlet{qblue}{blue!50!black}
\colorlet{lblue}{blue!50!black}
\colorlet{fblue}{blue!60!black}
\colorlet{fred}{red!60!black}
\DeclareGraphicsExtensions{.pdf,.png}


\usepackage{algorithm}
\usepackage[noend]{algpseudocode}
\usepackage{booktabs}
%{\renewcommand{\arraystretch}{1.2}
\newcommand{\comment}{\algorithmiccomment}
\DeclareMathAlphabet{\mathcal}{OMS}{cmsy}{m}{n}
\newcommand{\captionsize}{\normalfont}
\newcommand{\captionfont}{\captionsize}
%\usepackage{verbatimbox} %BB-AY this breaks my TeX setup
\usepackage{placeins}
% for fortran code listings

\pagestyle{fancy}
\renewcommand{\headrulewidth}{1pt}
\title{Accelerating the Computation of Geometric Constraints in Optimization}
\author{Benjamin J. Brelje}
\lhead{Benjamin J.~Brelje}
\chead{EECS587 Term Project}
\rhead{\today}
\begin{document}

\maketitle

\section{Introduction}
\qquad In aerospace engineering, particularly in commercial aviation, reducing cost and environmental impact is a major research focus.
My lab, the Multidisciplinary Design Optimization Lab (MDOLab), uses numerical optimization and computer simulation to create aircraft designs with minimum fuel consumption.
The output of one of our optimizations is an aircraft geometry file.

\qquad Previous graduate research resulted in efficient, massively parallel flow solvers and structural analysis tools.
These tools allow us to assess the weight and drag of a particular geometry.
The optimizer will run many of these aerostructural analyses until it converges on the optimal design.

\qquad Another important aspect of aircraft design is geometric constraints.
An example of a geometric constraint is that the aircraft must have enough volume in the wing for the mission's required fuel.
The focus of my research is making realistic geometric constraints easier to impose on our optimization problems.
In my previous work, I developed a Python package called \texttt{geograd} to compute the value of my new geometric constraints during optimization.
\texttt{Geograd} uses Tensorflow, a framework designed for machine learning, to compute my constraint in a pure SIMD fashion.
It was very fast when running on the GPU on my local machine, but we do not have access to GPU resources on our usual compute clusters.
The purpose of this term project is to greatly improve the performance of \texttt{geograd} by writing a brand-new, Fortran 90 implementation to take advantage of algorithm opportunities that cannot be posed in Tensorflow's SIMD language.

\qquad The first section will describe the specific computation in more detail and the current state of the software.
Second, I will describe a test case I constructed for my timings and benchmarking.
Following that, I will describe the development and verification of my baseline Fortran code.
Next, I will walk through each improvement I made to the algorithm and describe the dynamic load balancing problem that emerged during this stem.
Finally, I will describe the dynamic load balancing approach I used and present the final timings.

\qquad I was ultimately able to achieve over a 1350x speedup on my benchmark case compared to my original Tensorflow-based pure SIMD implementation. 

\section{Description of the Computation and Previous Work}
In numerical optimization, algorithms are designed to minimize some objective function $f(\textbf{x})$ by varying $\textbf{x}$ (a vector of design variables)
subject to $\textbf{g}(\textbf{x}) \leq 0$ (a vector of constraint values).
In aircraft optimization:
\begin{itemize}
 \item $f(\textbf{x})$ is generally fuel burn for a given mission
 \item $\textbf{x}$ are geometric parameters like wingspan, wing sweepback angle, and local thickness control points
 \item $\textbf{g}(\textbf{x})$ includes structural stress limits, stall speed requirements, and other requirements.
\end{itemize}
My research goal is to define a mathematical formulation for a component of $\textbf{g}$ which accounts for geometric constraints due to packing.
Notionally, my constraint $\textbf{g} \leq 0$ when the item I need to package fits entirely within the outer envelope of the airplane.

\qquad Our lab uses \emph{gradient-based} optimization because we have many (hundreds) of design variables $\textbf{x}$. 
Optimization algorithms which use gradient information are computationally cheaper than gradient-free methods on large-scale problems like ours.
When I say \emph{gradient information}, I am referring to the derivatives of the objective and constraints with respect to the design variables, namely:

\begin{equation}
    \frac{df(\textbf{x})}{d\textbf{x}}
\end{equation}

\begin{equation}
    \frac{d\textbf{g}(\textbf{x})}{d\textbf{x}}
\end{equation}

\qquad While gradient-based methods require orders of magnitude fewer function evaluations on large-scale problems, they require the researcher to define computational methods which include both function values and derivatives.
They also introduce requirements on the general behavior and smoothness of the functions.
Obviously, $f$ and $\textbf{g}$ must be once-differentiable.
Gradient-based methods perform best when $f$ and $\textbf{g}$ are deterministic.
Ideally, $f$ and $\textbf{g}$ are also $C^1$ continuous (smooth values and first derivatives), though adequate results can also be obtained with $C^0$ continuous functions (smooth values only) in some cases.

\qquad During the course of my research, I identified a mathematical formulation for a geometric constraint which meets these requirements.
We begin with triangulated representations of aircraft outer surface $A$ and an inner object $B$ to pack inside it, as pictured in Figure~\ref{fig:schematic}. 
$A$ and $B$ are each represented as lists of vertices $V_0$, $V_1$, $V_2$ of dimension 3 by $m$ or $n$ (the size of each mesh).

\begin{figure}[ht]
  \centering
  \includegraphics[width=0.45\linewidth]{figures/schematic}
  % \vspace{-30pt}
  \caption{A section view of an object $B$ inside wing $A$ and the minimum distance $d_\text{min}$ between them}
  \label{fig:schematic}
\end{figure}

When one object encloses another, the minimum distance $d_\text{min}$ between them is greater than zero by some margin.
Therefore, a first approach to a geometric constraint might involve computing $d_\text{min,AB} \geq 0 + \text{tol}$.
We can compute $d_\text{min}$ between two triangulated surfaces by computing the minimum distance between each individual pair of triangles.

\begin{equation}
  \label{eq:dmin}
  d_\text{min,AB} = \text{min}(d_{\text{min},ij}) \text{ for each triangle } i, \, j \text{ in surfaces } A, \, B
\end{equation}

The minimum distance between a pair of triangles can be found by a total of fifteen primitive tests between the vertices: six point-triangle tests, and nine line-line tests. % TODO cite ericson

\begin{equation}
  d_{\text{min},ij} = \text{min}(d_{\text{min},k}) \text{ for primitive test }k\text{1 to 15}
\end{equation}

The distances are then aggregated into a single scalar and constrained for feasibility using the alternate form of the Kresselmeier-Steinhauser (KS) function ~\ref{ks_geom}.
\begin{equation}
  \label{ks_geom}
  KS_\text{geom}(\textbf{x}) = \frac{1}{\rho} \textrm{ln} \Bigg[\sum_{j=1}^{J} e^{\rho\big(d_\text{min}\:(\textbf{x})-d_j\:(\textbf{x})\big)}\Bigg] - d_\text{min}\:(\textbf{x}) \leq 0 ,
\end{equation}

where $\textbf{x}$ is a vector of design variables,
$d_j(\textbf{x})$ is the $j\textrm{th}$ pairwise distance between $A$ and $B$,
$d_{min}(\textbf{x})$ is the minimum distance between $A$ and $B$, and
$\rho$ is a user-controllable constant.

\end{document}
